% ============================================================================
% WVC 2026 (Workshop de Visão Computacional) — UNICAMP submission.
% This file starts from the official IEEE conference template found by the
% author (version 062824, ieee.org) and keeps its preamble byte-for-byte:
% same core packages (cite; amsmath,amssymb,amsfonts; algorithmic; graphicx;
% textcomp; xcolor), same \IEEEoverridecommandlockouts and no a4paper option.
% The standard algorithm and url packages support the cascade float and
% the anonymous artifact link.
%
% Two deliberate content adaptations from our working draft
% (paper/wvc/wvc2026.tex) to fit this exact preamble:
%   - Tables converted from booktabs (\toprule/\midrule/\bottomrule) to
%     plain \hline rules, since booktabs is not in this template.
%   - Algorithm 1 converted from algpseudocode syntax (\State/\If/\ElsIf/
%     \Return) to plain algorithmic.sty syntax (\STATE/\IF/\ELSIF, with
%     "\STATE \textbf{return} ..." standing in for \Return, which plain
%     algorithmic.sty does not define).
%
% DOUBLE-BLIND NOTICE: author names/affiliations below are placeholders.
% Do not restore real identities until the camera-ready version after
% acceptance, since the CFP explicitly requires no identifying information
% in the blind submission, including in acknowledgments or linked material.
%
% NOTE: this paper is deliberately structured around an analysis question
% (does the necessity of audio-visual fusion vary across deployment
% conditions?) rather than around the system itself, and does not cite or
% reference the companion CONICT extended-abstract report on the same
% system per the author's explicit instruction. Keep the System Overview
% section compact; the stratified analysis in Results is this paper's
% primary contribution.
% ============================================================================
\documentclass[conference]{IEEEtran}
\IEEEoverridecommandlockouts
% The preceding line is only needed to identify funding in the first footnote. If that is unneeded, please comment it out.
%Template version as of 6/27/2024

\usepackage{cite}
\usepackage{amsmath,amssymb,amsfonts}
\usepackage{algorithm}
\usepackage{algorithmic}
\usepackage{graphicx}
\usepackage{textcomp}
\usepackage{xcolor}
\usepackage{url}
\def\BibTeX{{\rm B\kern-.05em{\sc i\kern-.025em b}\kern-.08em
    T\kern-.1667em\lower.7ex\hbox{E}\kern-.125emX}}
\begin{document}

\title{When Is Voice Not Enough? A Deployment-Scale Analysis of\\
Multimodal Necessity in Real-Time Speaker Attribution}

% TODO (camera-ready only): restore real author list / affiliation / emails.
% Blind-submission placeholder below.
\author{\IEEEauthorblockN{Anonymous Author(s)}
\IEEEauthorblockA{Affiliation withheld for double-blind review}
}

\maketitle

\begin{abstract}
Multimodal speaker attribution systems are typically evaluated with a
single aggregate accuracy or coverage figure, which implicitly assumes
that the benefit of fusing voice, face, and active-speaker cues is
uniform across deployment conditions. We test this assumption using
deployment logs from a real-time meeting-tracking system that
attributes each speech segment through an explicit, priority-ordered
audio-visual fusion cascade rather than a single learned model. Across
92 meeting sessions and 1,660 speech segments
collected over two months, we find that reliance on voice biometrics
varies from 5\% to 40\% of segments depending on expected group size,
while visual and session evidence extends coverage to every accepted
segment. The cascade sustains real-time throughput (89.6\% of segments
faster than real time). In a human-reviewed subset, its risk flag
separates unflagged named decisions (81.8\% exact agreement) from
flagged ones (24.8\%; Fisher's exact $p<0.001$), effectively
prioritizing review.
\end{abstract}

\begin{IEEEkeywords}
biometrics, speaker recognition, feature extraction, feature matching and inference, active speaker detection, real-time computer vision, audio-visual fusion
\end{IEEEkeywords}

\section{Introduction}

Multimodal speaker attribution combining voice biometrics, face
re-identification, and active speaker detection (ASD) is generally
reported to outperform any single modality
\cite{chung2017lipsync,tao2021talknet,liao2023lightasd}. Almost all
such evaluations, however, report one aggregate figure per system,
computed over a benchmark whose composition (how many people are in
frame, how long the deployment has been running, how well individual
speakers are enrolled) is fixed by construction. A meeting-transcription
system deployed in practice sees none of that stability. Group size
varies from a one-on-one call to a full room, and the same system may
run for months as its voice bank grows and its thresholds get tuned. Whether the reported benefit of fusion holds uniformly
across these conditions, or is concentrated in some and nearly absent
in others, is a question that a single aggregate figure cannot answer.

We study this question using deployment logs from a real-time,
single-machine multimodal meeting tracker that attributes each speech
segment through an explicit \emph{priority cascade}, an ordered
sequence of decision rules in which the first rule whose evidence
requirement is met determines the final identity, rather than a single
learned fusion model. At its core, the underlying problem is a
multimodal \emph{biometric feature matching and inference} task, in
which voice and face embeddings are extracted independently by two
different feature-extraction pipelines (ECAPA-TDNN acoustic embeddings
and EdgeFace visual embeddings) and must be matched, cross-checked, and
reconciled online against a per-session identity registry. We use the
cascade's logged decisions, collected over 92 real, unscripted sessions
and 1{,}660 speech segments spanning just over two months, as an
instrument to measure how the necessity of each modality shifts across
deployment conditions rather than to report one more aggregate accuracy
number.

Our contributions are the following.

\begin{enumerate}
\item Evidence that reliance on voice biometrics is highly non-uniform
across group sizes, ranging from 5\% to 40\% of segments, with a
structural explanation tied to which cascade rule takes precedence at
each group size, rather than a smooth trend.
\item A demonstration that an apparent decline in voice-led
resolution over a multi-month deployment is substantially confounded
by a shift in session composition, a methodological caution for
interpreting raw temporal trends in deployment logs without
controlling for this kind of shift.
\item A statistically tested human validation of the lightweight
self-diagnostic risk monitor, together with a real-time performance
characterization of the cascade under deployment conditions.
\item A reproducible evaluation protocol, with anonymized segment-level
tables and scripts that regenerate the reported metrics.
\end{enumerate}

\section{Related Work}

\textbf{Speaker diarization and verification.} pyannote.audio
\cite{bredin2020pyannote,bredin2023pyannote31} provides the diarization
backbone used in our pipeline, segmenting and clustering speech into
per-chunk speaker labels; diarization alone does not solve identity
persistence across non-adjacent processing windows. ECAPA-TDNN
\cite{desplanques2020ecapa}, trained on VoxCeleb, is the standard
architecture for the fixed-dimensional speaker embeddings we use both
for enrolled-speaker matching and for a persistent within-session
tracker for unenrolled speakers.

\textbf{Active speaker detection.} SyncNet-style lip-sync models
\cite{chung2017lipsync} and TalkNet \cite{tao2021talknet} established
the audio-visual correlation approach to ASD; Light-ASD
\cite{liao2023lightasd} is a lightweight CVPR 2023 model designed for
real-time inference, which we use as the primary ASD backend. We use
YOLOv8 \cite{jocher2023yolov8} for face detection, ByteTrack
\cite{zhang2022bytetrack} for short-term tracking, and EdgeFace
\cite{george2024edgeface} for lightweight face-recognition embeddings.

\textbf{Audio-visual fusion for speaker identity.} Most prior work on
audio-visual speaker identification targets offline, dataset-level
fusion, such as re-scoring diarization with ASD or jointly training
audio-visual embeddings, and reports a single accuracy figure over a
benchmark of fixed composition. Our setting differs in requiring an
\emph{online}, segment-by-segment decision under partial observability,
where a face, a voice match, or both may be missing or unreliable for
any given segment, which led us to a rule-based, explicitly ordered
fusion strategy rather than a single learned fusion score.

\textbf{Generalization across deployment conditions.} Benchmarks such
as AVA-ActiveSpeaker \cite{roth2020ava} and VoxConverse
\cite{chung2020voxconverse} fix group size and recording conditions by
construction, so a system's reported accuracy is, by design, an average
over whatever mixture of conditions the benchmark happens to contain.
We are not aware of prior work that stratifies a deployed multimodal
attribution system's behavior by group size or by deployment time to
test whether a single aggregate figure is representative, which is the
gap this paper addresses.

\section{System Overview}
\label{sec:system}

The system runs two cooperating threads sharing an in-memory state
object. A \emph{video thread} performs face detection (YOLOv8),
tracking (ByteTrack), face re-identification (EdgeFace), and per-frame
ASD updates. An \emph{audio thread} performs voice activity detection,
diarization-based session hints (pyannote.audio), transcription
(faster-whisper \cite{radford2022whisper}, Portuguese),
and speaker verification (ECAPA-TDNN) on each detected speech segment,
querying the video thread's state for currently visible faces and ASD
results. When diarization indicates more than one concurrent speaker,
the system applies SepFormer \cite{subakan2021sepformer} speech
separation before re-verifying each separated stream.

Algorithm~\ref{alg:cascade} summarizes the fusion cascade that
attributes each segment. Each rule has a different evidence
requirement, and the cascade stops at the first rule whose precondition
is satisfied. It prefers voice biometrics whenever the match is
confident enough, since a voice match is invariant to camera framing
and remains available for off-camera speakers, and falls back to
progressively weaker, more circumstantial evidence otherwise. Two
cross-cutting passes run afterward. A \textbf{gender-consistency check}
rejects or reroutes visually derived decisions that contradict an
unambiguous audio-derived gender signal, skipped when voice-biometric
confidence is very high ($\geq 0.90$). A \textbf{turn-merging} pass
merges consecutive same-diarization-label segments only if their
voice-embedding similarity exceeds 0.40, since diarization occasionally
mislabels two different speakers under the same local cluster.

The two accepted-match tiers in Algorithm~\ref{alg:cascade} are
separated by the confirmation threshold $\tau_c$, which is configured
independently of the acceptance threshold $\tau$. The deployment ran at
$\tau = 0.80$, $\tau_c = 0.70$ for 97\% of segments, and at that
operating point every accepted match also clears $\tau_c$, so the two
tiers collapse into one and \textsc{Verifier-Mod} is reached only in a
small set of sessions configured with a looser $\tau = 0.55$. Because
this split is an artifact of configuration rather than of evidence
strength, we report all three verifier tiers together as a single
voice-led decision family throughout \S\ref{sec:results}.

\begin{algorithm}[h]
\caption{Speaker attribution cascade (simplified)}
\label{alg:cascade}
\begin{algorithmic}[1]
\STATE $v \gets$ accepted ECAPA match with confidence $c$; accepted only
  if $c \geq \tau$ and its margin over the runner-up is $\geq 0.07$
\STATE $r \gets$ best raw ECAPA candidate (name, confidence), accepted or not
\STATE $g \gets$ audio-derived gender estimate
\STATE $a \gets$ ASD query over [segment start $-$ 0.4s, end $+$ 0.4s]
\IF{$v$ is not \textsc{None}, $c \geq \tau_c$, \textbf{and} gender-consistent}
  \STATE \textbf{return} \textsc{Verifier-High} ($v$.name)
\ELSIF{$v$ is not \textsc{None} \textbf{and} gender-consistent}
  \STATE \textbf{return} \textsc{Verifier-Mod} ($v$.name)
\ELSIF{$r$.conf $\geq 0.65$ \textbf{and} (ASD or gender corroborates)}
  \STATE \textbf{return} \textsc{Verifier-Weak} ($r$.name)
\ELSIF{$a$ is not \textsc{None}}
  \STATE \textbf{return} \textsc{ASD} (face bound to $a$)
\ELSIF{exactly 1 face visible, known identity, $\leq$1 expected speaker}
  \STATE \textbf{return} \textsc{Face-Only}
\ELSIF{exactly 1 face visible, $r$.conf $\geq 0.35$, $\leq$1 expected speaker}
  \STATE \textbf{return} \textsc{Single-Face} (heuristic)
\ELSIF{ASD best-guess is dominant ($\geq$1.5$\times$ runner-up)}
  \STATE \textbf{return} \textsc{ASD-Guess}
\ELSIF{diarization session hint bound to a known face}
  \STATE \textbf{return} \textsc{Session-Face}
\ELSIF{diarization session hint available}
  \STATE \textbf{return} \textsc{Session-Hint} (anonymous)
\ELSE
  \STATE \textbf{return} \textsc{Session-Tracker} (persistent voice cluster)
\ENDIF
\STATE \textit{(Gender cross-check / reroute applied as a final pass)}
\end{algorithmic}
\end{algorithm}

ASD itself uses a two-backend design, with Light-ASD
\cite{liao2023lightasd} as the primary backend and a model-free
pixel-difference fallback for windows in which a neural score is not
yet available. This design preserves continuous online operation. The
fallback also rejects occlusion-related motion by requiring the mouth
region to move at least 1.8$\times$ more than the upper face before the
motion is treated as speech-related.

\section{Experimental Setup}
\label{sec:setup}

The evaluation uses an internal development dataset collected during
real, unscripted operation from 16 March to 22 May 2026. A session was
included when it contained at least one accepted speech segment. This
criterion retained 92 sessions and 1{,}660 segments; 14 empty sessions
were excluded before analysis. The sessions were recorded in
Portuguese with consumer webcam and microphone hardware under varying
lighting, framing, overlap, and enrollment availability. Their
configured group-size distribution was 19 one-speaker, 7 two-speaker,
4 three-speaker, and 62 four-speaker sessions. Processing used one
machine with an Intel Core i5-11400H, 8\,GB RAM, and an NVIDIA RTX 3060
Laptop GPU with 6\,GB VRAM.

Audio is captured at 16\,kHz. The first 8\,s of a session form the
initial processing window; subsequent windows contain 15\,s and advance
by 5\,s, preserving 10\,s of context while transcribing only the new
region. pyannote.audio 3.1 receives the configured number of speakers as
both its minimum and maximum, uses clustering threshold 0.60 and
segmentation threshold 0.40, and retains turns of at least 0.5\,s. The
transcription stage uses beam size 5, speech VAD, a 500\,ms minimum
silence duration, and 300\,ms speech padding. The model configured for
each retained session was recorded in the database: 87 sessions used
\texttt{large-v3-turbo}, 3 used \texttt{large-v3}, and 2 used
\texttt{medium}. This distribution is preserved in the analysis rather
than treating the deployment as a single-model benchmark.

\begin{table}[h]
\caption{Principal operating configuration.}
\begin{center}
\footnotesize
\begin{tabular}{ll}
\hline
\textbf{Stage} & \textbf{Configuration} \\
\hline
Audio & 16\,kHz; 8\,s initial; 15\,s window; 5\,s step \\
Diarization & pyannote 3.1; 0.60 clustering; 0.40 segmentation \\
Transcription & faster-whisper; beam 5; VAD enabled \\
Voice match & $\tau=0.80$ (88 sess.); $0.55$ (4); margin 0.07 \\
Confirmation & $\tau_c=0.70$ in all 92 sessions \\
ASD query & speech interval expanded by $\pm0.4$\,s \\
\hline
\end{tabular}
\label{tab:configuration}
\end{center}
\end{table}

For every accepted segment, SQLite stores the session identifier,
segment index, final decision type, identity, raw voice confidence,
ASD and face indicators, risk flag, speech duration, and transcription,
verification, and total processing times. We group decisions into five
families: voice-only, voice with corroboration, active-speaker, face or
gender, and session context. This segment-level provenance makes each
aggregate in the results traceable to the rule that produced it.

The five evidence families are mutually exclusive and exhaustive.
\textsc{Verifier-High} and \textsc{Verifier-Mod} form the voice-only
family; \textsc{Verifier-Weak} forms voice with corroboration;
\textsc{ASD} and \textsc{ASD-Guess} form active-speaker evidence;
face and gender decisions form the visual identity family; and
\textsc{Session-Face}, \textsc{Session-Hint}, and
\textsc{Session-Tracker} form session context. This mapping prevents a
segment from contributing to more than one percentage in the aggregate
coverage analysis.

For segment $i$, with total processing time $t_i$ in milliseconds and
speech duration $d_i$ in seconds, the real-time factor is
\begin{equation}
    \mathrm{RTF}_i=\frac{t_i}{1000d_i}.
\end{equation}
An RTF below 1 means that processing completed faster than the speech
segment itself. We report the median RTF and the proportion of segments
with $\mathrm{RTF}<1$, together with stage-level medians.

Identity agreement was evaluated on 330 manually reviewed segments
from 24 sessions. When more than one review existed for a session, only
the latest was retained. Exact agreement requires equality between the
predicted and corrected identity after whitespace and case
normalization. The named subset excludes anonymous diarization labels.
Proportions are accompanied by 95\% Wilson intervals
\cite{wilson1927}; the association between the risk flag and exact
agreement is tested with a two-sided Fisher exact test at
$\alpha=0.05$.

DER and WER require time-aligned speaker turns and reference word
transcripts, respectively \cite{bredin2017metrics}. Those annotations
were not part of the deployment log, and the fusion cascade neither
changes the Whisper words nor emits a reference-aligned diarization
timeline. We therefore report metrics aligned with the evaluated task:
decision provenance, exact named-identity agreement, and RTF, without
presenting component metrics as end-to-end attribution scores.

\begin{figure}[h]
\centering
\fbox{\parbox{0.92\columnwidth}{\centering\footnotesize
Audio/video capture $\rightarrow$ diarization, transcription, face and
ASD processing $\rightarrow$ deterministic fusion cascade $\rightarrow$
segment-level SQLite log $\rightarrow$ grouped metrics, confidence
intervals, and human-validation test}}
\caption{Evaluation flow from multimodal capture to reproducible metrics.}
\label{fig:evaluation-flow}
\end{figure}

The anonymous artifact at
\url{https://github.com/anonimo826/anonimo} contains the system snapshot,
parameter file, data dictionary, privacy-preserving derived tables, and
an analysis script that regenerates the counts, Wilson intervals, RTF
statistics, and Fisher test with one command. The released tables use
sequential session codes and binary identity-agreement fields; they do
not contain audio, video, names, transcripts, absolute dates, or
biometric embeddings. SHA-256 hashes and a manifest allow researchers
to verify that the analyzed artifacts remain unchanged.

\section{Results}
\label{sec:results}

\subsection{Voice-led resolution is non-monotonic in group size}
\label{sec:bysize}

Across the whole deployment, voice evidence leads 8.6\% of the 1{,}660
decisions: 4.0\% are accepted voice matches and 4.6\% are voice matches
corroborated by ASD or gender. The remaining 91.4\% demonstrate the
coverage gained from ASD, face re-identification, gender consistency,
and persistent session context. Table~\ref{tab:byspeakers} stratifies
the same segments by the configured group size. The
voice-led rate is not a smooth function of group size, and each value
has a structural explanation rooted in the cascade itself
(Algorithm~\ref{alg:cascade}), not merely in how hard voice matching
happens to be. Single-speaker sessions show a low voice-led rate
(11.3\%, 95\% CI [5.8, 20.7]) because the cascade's face-identity rules are unlocked
precisely when at most one speaker is expected, so voice matching is
rarely the deciding rule even when it would have agreed. Two-speaker
sessions show the highest voice-led rate observed (39.8\%, 95\% CI
[31.5, 48.8]), plausibly
because a two-person voice bank is the easiest case to disambiguate and
these sessions had the highest proportion of fully enrolled
participants. Three-speaker sessions lean heavily on the
gender-consistency layer (23.3\%) and, as Table~\ref{tab:fprisk} shows
below, also carry the highest self-flagged risk of any group size.
Four-speaker sessions, which contribute 82\% of all segments, drop to
the lowest voice-led rate (5.0\%, 95\% CI [3.9, 6.3]) while ASD carries
the largest share (36.8\%). Together, these results show that the
cascade adapts its evidence source to session composition and that the
visual branches are most active in the dominant four-speaker setting.

\begin{table}[h]
\caption{Decision source by number of expected speakers per session.
Columns report the three main decision families and do not sum to
100\%; the balance is carried by the face-only, session-hint and
session-tracker rules.}
\begin{center}
\begin{tabular}{lrrrrr}
\hline
\textbf{Speakers} & \textbf{Sess.} & \textbf{Segs.} & \textbf{Voice} & \textbf{ASD} & \textbf{Gender} \\
\hline
1 & 19 &   71 & 11.3\% & 16.9\% &  1.4\% \\
2 &  7 &  118 & 39.8\% & 35.6\% &  0.8\% \\
3 &  4 &  103 & 19.4\% & 22.3\% & 23.3\% \\
4 & 62 & 1368 &  5.0\% & 36.8\% &  6.9\% \\
\hline
\end{tabular}
\label{tab:byspeakers}
\end{center}
\end{table}

Table~\ref{tab:fprisk} cross-checks this with the self-diagnostic risk
flag (\S\ref{sec:fprisk}). The risk rate is also non-monotonic and
peaks at three speakers (68.9\%), well above both its
neighbors, matching the group size where the gender-safety layer had
to work hardest in Table~\ref{tab:byspeakers}. Three-speaker meetings
therefore receive the highest review priority from the monitor. This
behavior is consistent with the strong participation of the
gender-consistency layer in the same group.

\begin{table}[h]
\caption{Self-flagged review-priority rate by expected speakers.}
\begin{center}
\begin{tabular}{lrrr}
\hline
\textbf{Speakers} & \textbf{Sess.} & \textbf{Segments} & \textbf{Flagged \%} \\
\hline
1 & 19 &   71 & 23.9\% \\
2 &  7 &  118 &  5.9\% \\
3 &  4 &  103 & 68.9\% \\
4 & 62 & 1368 & 28.1\% \\
\hline
\end{tabular}
\label{tab:fprisk}
\end{center}
\end{table}

\subsection{Apparent temporal decline is a composition effect}
\label{sec:bytime}

A naive reading of the deployment log by calendar week shows what
looks like a steep decline in voice-led resolution. Voice-led
decisions accounted for 37--42\% of segments in the first two weeks,
falling to 0--2\% in the final three weeks recorded. Read on its own,
this trend invites a misleading story, namely that voice matching
degraded as the deployment progressed.

Breaking down the same weeks by session composition rules this out.
The first two weeks were composed entirely of one- and two-speaker
sessions (100\% in each week), the group sizes Table~\ref{tab:byspeakers}
identifies as having the highest voice-led rates. From the third week
onward, the deployment shifted to sessions with three or four expected
speakers, reaching 90\% four-speaker segments in the third week and
staying above 80\% in every subsequent week with more than 100
segments. The decline in voice-led resolution over time is, to a
substantial degree, the same non-monotonic group-size effect from
\S\ref{sec:bysize} appearing again under a different label, because the
deployment's session mix changed, not because voice biometrics got
worse at its job.

The self-flagged risk rate over the same weeks makes the point a
second time, and in the opposite direction. It was zero in the first
two weeks, rose to 39\% and then 52\% as multi-speaker sessions came to
dominate, and fell back to 14\% in the final week. Read as a time
series alone, the rise and the fall would each invite a story about the
system changing. Both are largely the composition shift again, with the
final-week drop additionally consistent with threshold and gender-check
refinements made late in the period, a genuine improvement running in
parallel with, and partly obscured by, that shift. We take this as a
methodological caution. A raw trend line
over a real deployment log can be confounded by what the deployment
happened to contain in each period, and stratifying by a plausible
confounder, here group size, before drawing conclusions from time alone
is necessary, not optional.

\subsection{Real-time performance}

The cascade processed 1{,}488 of 1{,}660 segments (89.6\%) faster than
real time, with median RTF 0.24. Median total processing time was
633.1\,ms per segment; transcription accounted for 513.6\,ms and voice
verification for 29.7\,ms. These measurements show that multimodal
attribution remains compatible with live meeting operation on a single
consumer laptop GPU.

\subsection{ASD backend accuracy against the published literature}
\label{sec:columbia}

The ASD backend is used exactly as released, with no retraining or
fine-tuning on our data, so the figures its authors report describe the
same weights that run in our pipeline. The pretrained Light-ASD
checkpoint we use is reported by its authors to
reach 94.06\% mAP on the AVA-ActiveSpeaker validation set and an
average F1 of 81.1\% on the Columbia ASD dataset
\cite{chakravarty2016columbia}, a small, five-panelist benchmark widely
used for out-of-domain ASD evaluation \cite{liao2023lightasd}. Both are
fixed-composition benchmarks of the kind this paper's stratified
deployment analysis is meant to complement, and neither speaks to how
often ASD is the rule that decides an identity, which is what
Table~\ref{tab:byspeakers} measures.

\subsection{Statistical validation of the review-priority monitor}
\label{sec:fprisk}

Independently of the cascade, the system computes a rule-based risk
flag for every segment that received a named identity. Four factors
raise it. A decision of a structurally weak type
(\textsc{Verifier-Weak}, \textsc{Face-Only}, \textsc{Single-Face})
raises its review priority, as does the number of distinct identified speakers
in the session exceeding the configured expected count. A name drawn
from the enrolled voice bank whose raw voice confidence never reached
the acceptance threshold $\tau$ receives high priority, as does a
contradiction between the attributed name's gender and the
audio-derived gender estimate. The monitor is a deterministic function
of decision metadata and therefore operates without labels. Across the
deployment, 479 of 1{,}660 segments (28.9\%) were flagged for priority
review.

The human-reviewed subset confirms that this prioritization is
informative. Among named decisions, 9 of 11 unflagged cases agreed with
the human correction (81.8\%, Wilson 95\% CI [52.3, 94.9]), compared
with 28 of 113 flagged cases (24.8\%, 95\% CI [17.7, 33.5]). Fisher's
exact test rejects independence between the flag and exact agreement
($p=2.94\times10^{-4}$). Thus, the monitor identifies a substantially
different review-priority group and provides a practical way to focus
human verification effort.

\section{Discussion}

The results support a condition-aware view of multimodal attribution.
The same cascade moves naturally between voice, ASD, face, and session
context as group composition changes. This is particularly valuable in
four-speaker sessions, which represent 82.4\% of the evaluated
segments: visual or session evidence participates in 95.0\% of their
decisions, preserving coverage when voice alone is not decisive.

The combination of traceable decision rules, a statistically validated
review-priority monitor, and real-time execution offers three practical
benefits. Each identity can be audited back to its evidence source,
segments can be ranked for efficient human review, and the complete
analysis can be regenerated from privacy-preserving tables. These
properties make the system a reproducible foundation for future
condition-aware threshold selection and larger annotated evaluations.

\section{Conclusion}

In 92 real meeting sessions and 1{,}660 speech segments, the system
combined complementary modalities while sustaining real-time operation
for 89.6\% of segments. Voice-led decisions ranged from 5\% to 40\%
across group sizes, demonstrating why deployment composition should be
reported together with aggregate coverage. The review-priority monitor
also separated named decisions with 81.8\% and 24.8\% exact agreement
($p<0.001$), turning decision metadata into an effective aid for human
verification. Overall, the deterministic cascade provides broad
coverage, traceable decisions, real-time execution, and a reproducible
evaluation path for multimodal meeting intelligence.

% ---------------------------------------------------------------------
% CAMERA-READY ONLY. Do NOT uncomment before acceptance: the WVC CFP
% requires that no identifying information appear in the blind
% submission, and this text names the collaborating institutions.
% Adapted from the Portuguese "Agradecimentos" of the companion CONICT
% submission (paper/conict/main.tex), which must stay anonymous for the
% same reason until its own final version.
%
% Replace the paragraph above with:
%
% The authors thank the CEPIN-CIS research center, the IT University of
% Copenhagen, and the researchers at GN Group for making this study
% possible, and for the technical support and resources that were
% essential to carrying it out. We also thank everyone who contributed,
% directly or indirectly, to this work.
%
% NOTE: write out "IT University of Copenhagen" rather than the bare
% acronym ITU, which an international audience will read as the
% International Telecommunication Union.
% ---------------------------------------------------------------------

\bibliographystyle{IEEEtran}
\bibliography{references}

\end{document}
